\documentclass[12pt]{article}
\usepackage{preamble}
\usepackage{longtable}

\geometry{a4paper, left=2.5cm, right=2.5cm, top=2cm, bottom=2.5cm}

\pagestyle{plain}

\newcommand{\placeholder}[1]{{\color{magenta}#1}}

\chead{
    \begin{minipage}{1\linewidth}
        \begin{wrapfigure}{r}{0pt}
            \includegraphics[height=1cm]{images/logo}
        \end{wrapfigure}
        {
            \centering
            \sffamily\scriptsize
            \textbf{
                Санкт-Петербургский национальный исследовательский университет \\
                информационных технологий, механики и оптики}
            %th3p4g
            \vspace{2mm}

            \quad\quad\quad\quad\quad\quad\ \textbf{УЧЕБНЫЙ ЦЕНТР ОБЩЕЙ ФИЗИКИ ФТФ}
        }
    \end{minipage}
}



\begin{document}
    \vspace*{2\baselineskip}

    \thispagestyle{fancy}

    \noindent
    \textbf{Группа} \underline{M3200\hspace{4.5cm}} \hfill \textbf{К работе допущен} \underline{\hspace{4cm}} \\[0.5cm]
    \textbf{Студенты} \underline{Комашко А., Семёнов Д.\hspace{0.45cm}} \hfill \textbf{Работа выполнена} \underline{\hspace{4cm}} \\[0.5cm]
    \textbf{Преподаватель} \underline{Шоев В. И.\hspace{1.9cm}} \hfill \textbf{Отчет принят} \underline{\hspace{4.85cm}} \\



    \begin{center}
    {\huge \textbf{Рабочий протокол и отчёт по\\ лабораторной работе № 3.03}}

        \smallvspace

        {\Large Определение удельного заряда электрона}
    \end{center}


    \noindent
    1. \textbf{Цель работы.}

    \begin{enumerate}
        \item \placeholder{Цель №1}

        \item \placeholder{Цель №N}
    \end{enumerate}

    \mediumvspace

    \noindent
    2. \textbf{Задачи, решаемые при выполнении работы.}

    \begin{enumerate} 
        \item Провести измерения зависимости анодного тока $I_a$ вакуумного
        диода от величины тока в соленоиде при различных значениях
        анодного напряжения.
        \item Найти значение коэффициента связи между током соленоида
        и магнитным полем $B$ внутри него.
        \item Построить графики зависимостей $I_a$ от $B$ и определить по
        ним величины критических полей для каждого значения анодного
        напряжения.
        \item По значениям критического поля найти величину удельного
        заряда электрона и оценить ее погрешность.
    \end{enumerate}

    \mediumvspace

    \noindent
    3. \textbf{Объект исследования: } 
    
    Поток электронов внутри диода под воздействием магнитного поля соленоида

    \mediumvspace

    \noindent
    4. \textbf{Метод экспериментального исследования.}

    \placeholder{Метод исследования}

    \mediumvspace

    \noindent
    5. \textbf{Рабочие формулы и исходные данные.}

    \placeholder{Рабочая формула №1:
\[
\vec{F} = m\vec{a}
\]

Рабочая формула №2:
\[
E = mc^2
\]}

    \mediumvspace

    \noindent
    6. \textbf{Измерительные приборы.}

    \smallvspace

    \begin{center}
    \begin{tabular}{|c|m{5cm}|C{2cm}|C{2cm}|C{2cm}|c|}
        \hline
        № п/п & Наименование            & Предел\newlineизмерений & Цена\newlineделения     & $\Delta_n$ \\
        \hline
        1     & Мультиметр MY64 & 0 - 10  A     & 1  мA  & 1  мА     \\
        \hline
        2     & Мультиметр MY65 & 0 - 2  мА     & 1 мкА  & 1 мкА    \\
        \hline

    \end{tabular}

    \smallvspace

    \textit{Таблица 1.} Измерительные приборы
\end{center}

    \mediumvspace

    \noindent
    7. \textbf{Схема установки. (перечень схем, которые составляют Приложение 1).}

    \hyperlink{schema1}{Схема установки} прилагается в Приложении 1

    \mediumvspace

    \noindent
    8. \textbf{Результаты прямых измерений и их обработки (таблицы, примеры расчетов).}

    \begin{itemize}
    \item Прилагается \hyperlink{table2}{Таблица 2} в Приложении 2
\end{itemize}


    \mediumvspace

    \noindent
    9. \textbf{Расчет результатов косвенных измерений (таблицы, примеры расчетов).}

    \begin{itemize}
    \item \placeholder{Интересные вычисления: $E = 0.4 \cdot (3 \cdot 10^8)^2 = 3.6 \cdot 10^{16}$}

    \item Прилагается \hyperlink{tableM}{Таблица M} в Приложении 2
\end{itemize}

    \mediumvspace

    \noindent
    10. \textbf{Расчёт погрешности измерений.}

    \mediumvspace

    \noindent
    11. \textbf{Графики (перечень графиков, которые составляют Приложение 3).}

    \hyperlink{diagram1}{График 1} прилагаются в Приложении 3

    \mediumvspace

    \noindent
    12. \textbf{Окончательные результаты.}

    $\frac{e}{m} = (1.66 \pm 0.28) \cdot 10^{11} $ Кл/кг $\qquad \varepsilon_{\frac{e}{m}} = 16.7\% \qquad \alpha = 0.95$

    \mediumvspace

    \noindent
    13. \textbf{Выводы и анализ результатов работы.}

    \placeholder{На основе результатов вышло пу-пу-пу}

    \clearpage

    \begin{center}
        \LARGE
        \textbf{Приложение 1. Схема установки}
    \end{center}

    \mediumvspace

    \hypertarget{schema1}{}

\begin{center}
    \includegraphics[width=10cm]{images/logo}

    \smallvspace

    \textit{Рисунок 1.} \placeholder{Общий вид ИТМО}
\end{center}

\begin{enumerate}
    \item \placeholder{ИТМО 1}
    \item \placeholder{ИТМО 2}
    \item \placeholder{ИТМО 3}
\end{enumerate}






    \clearpage

    \begin{center}
        \LARGE
        \textbf{Приложение 2. Таблицы измерений и расчётов}
    \end{center}

    \mediumvspace

    \begin{center}
    \hypertarget{table2}{}

    \renewcommand{\arraystretch}{1.8}

    \begin{longtable}{|C{1.6cm}|C{2cm}|C{2cm}|C{2cm}|C{2cm}|C{2cm}|C{2cm}|}
        \hline
        \multirow{3}{*}{№ опыта} & \multicolumn{6}{c|}{Анодное напряжение} \\
        \cline{2-7}
        & \multicolumn{2}{c|}{$U = 9$ В} & \multicolumn{2}{c|}{$U = 10.5$ В} & \multicolumn{2}{c|}{$U = 13.5$ В} \\
        \cline{2-7}
        & $I_L$, мА & $I_a$, мА & $I_L$, мА & $I_a$, мА & $I_L$, мА & $I_a$, мА \\
        \hline
        1 & $0$ & $0.218$ & $0$ & $0.264$ & $0$ & $0.357$\\
        \hline
        2 & $23$ & $0.218$ & $29$ & $0.264$ & $24$ & $0.357$\\
        \hline
        3 & $42$ & $0.218$ & $59$ & $0.264$ & $52$ & $0.357$\\
        \hline
        4 & $67$ & $0.218$ & $80$ & $0.264$ & $77$ & $0.358$\\
        \hline
        5 & $95$ & $0.219$ & $102$ & $0.264$ & $97$ & $0.358$\\
        \hline
        6 & $123$ & $0.219$ & $131$ & $0.265$ & $118$ & $0.358$\\
        \hline
        7 & $145$ & $0.218$ & $155$ & $0.266$ & $148$ & $0.359$\\
        \hline
        8 & $165$ & $0.217$ & $184$ & $0.264$ & $176$ & $0.360$\\
        \hline
        9 & $192$ & $0.212$ & $209$ & $0.257$ & $202$ & $0.355$\\
        \hline
        10 & $213$ & $0.203$ & $230$ & $0.228$ & $231$ & $0.339$\\
        \hline
        11 & $239$ & $0.153$ & $259$ & $0.162$ & $256$ & $0.254$\\
        \hline
        12 & $260$ & $0.126$ & $284$ & $0.125$ & $279$ & $0.198$\\
        \hline
        13 & $283$ & $0.098$ & $314$ & $0.103$ & $299$ & $0.167$\\
        \hline
        14 & $310$ & $0.081$ & $342$ & $0.088$ & $322$ & $0.148$\\
        \hline
        15 & $332$ & $0.070$ & $362$ & $0.079$ & $346$ & $0.134$\\
        \hline
        16 & $362$ & $0.058$ & $383$ & $0.070$ & $374$ & $0.114$\\
        \hline
        17 & $382$ & $0.052$ & $403$ & $0.063$ & $400$ & $0.100$\\
        \hline
        18 & $405$ & $0.046$ & $424$ & $0.056$ & $424$ & $0.091$\\
        \hline
        19 & $427$ & $0.041$ & $445$ & $0.051$ & $447$ & $0.082$\\
        \hline
        20 & $449$ & $0.037$ & $467$ & $0.047$ & $472$ & $0.073$\\
        \hline
    \end{longtable}

    \smallvspace

    \textit{Таблица 4.} Зависимость напряжения $U_R$ от тока в соленоиде
\end{center}

    \begin{center}
    \hypertarget{table3}{}

    \renewcommand{\arraystretch}{1.8}

    \begin{tabular}{|c|C{2.5cm}|C{2.5cm}|C{1.8cm}|}
        \hline
        $U$, В & $I_{L_{\text{кр.}}}$, мкА & $B_\text{кр.}$, мкТл & $e/m$, Кл/кг  \\
        \hline
        & & & \\
        \hline
        & & & \\
        \hline
        & & & \\
        \hline
    \end{tabular}

    \smallvspace

    \textit{Таблица \placeholder{M}.} \placeholder{Какие-то интересные расчёты}

\end{center}

    \clearpage

    \begin{center}
        \LARGE
        \textbf{Приложение 3. Графики}
    \end{center}

    \mediumvspace

    \hypertarget{diagram1}{}

\begin{center}
    \includegraphics[width=15cm]{images/logo}

    \smallvspace

    \textit{График 1.} \placeholder{Зависимость $ITMO = I(t) \cdot \frac{M}{O}$}
\end{center}



\end{document}